%!TEX root = thesis.tex
\chapter{Conclusions and Future Work}
\label{chap:conclusion}



This chapter presents the conclusions drawn from the study and discusses their implications for the management of large-scale time-series data generated by the Smart Ocean underwater sensor network. The study designed and implemented a benchmarking framework to evaluate the ingestion performance, query performance and storage efficiency of MongoDB, InfluxDB and TimescaleDB under representative environmental observation workloads. Based on the experimental results presented in Chapter 5, the chapter summarizes the major findings in relation to the research objectives and research questions of the study. It also discusses the limitations of the benchmarking approach and identifies areas where further investigation may be required. Finally, recommendations are provided for database selection in Smart Ocean environments and for future research aimed at extending the evaluation of time-series database technologies for environmental monitoring applications.


\section{Conclusions}
The goal of this study was to evaluate the suitability of selected database technologies for managing the large volumes of time-series data generated by the Smart Ocean underwater sensor network. As environmental monitoring systems continue to produce increasing quantities of timestamped sensor observations, selecting an appropriate database platform becomes an important consideration for ensuring efficient data storage, retrieval and long-term management. To address this challenge, the study designed and implemented a benchmarking framework capable of evaluating the performance of MongoDB, InfluxDB and TimescaleDB under representative environmental observation workloads.

A major outcome of the research was the successful development and implementation of a reproducible benchmarking framework that enabled consistent evaluation across the three database systems. The framework provided a controlled environment in which identical datasets could be ingested, queried and analyzed under equivalent experimental conditions. Through the use of standardized datasets, common performance metrics, containerized deployment and automated benchmark execution, the framework generated a reliable basis for comparing the behavior of the evaluated database technologies. The implementation demonstrated that it is possible to perform systematic benchmarking of heterogeneous database systems while maintaining consistency in workload execution and result collection.

The benchmark results revealed that database architecture had a significant influence on performance outcomes. Although all three database systems successfully stored and retrieved the benchmark datasets, substantial differences were observed in ingestion performance, query latency and storage efficiency. These differences became more pronounced as dataset sizes increased, indicating that the internal design of a database plays an important role in determining how effectively it scales under growing time-series workloads.


Among the evaluated systems, MongoDB consistently demonstrated the strongest performance across the metrics examined in this study. The database achieved the shortest execution times during data ingestion, the highest observation throughput and data throughput values, the lowest latest-value query latencies and the smallest storage footprint. These results indicate that MongoDB was highly effective at handling the environmental observation workloads used in the benchmark. Its document-oriented architecture enabled efficient storage and retrieval of the benchmark data and allowed it to maintain strong performance as dataset sizes increased. Within the scope of the evaluated workloads, MongoDB provided the most favourable overall combination of ingestion efficiency, query responsiveness and storage utilization.

InfluxDB demonstrated competitive and balanced performance throughout the experiments. While it did not achieve the ingestion throughput or storage efficiency observed in MongoDB, it maintained stable behavior across increasing dataset sizes and consistently produced lower query latencies than TimescaleDB. The benchmark results showed that InfluxDB scaled predictably as workload volume increased and maintained relatively consistent performance across repeated benchmark executions. These findings suggest that InfluxDB provides a balanced approach to managing time-series data and remains a viable option for environmental monitoring applications requiring dependable ingestion and retrieval performance.

TimescaleDB successfully supported all benchmark workloads and correctly executed both ingestion and query operations. However, under the conditions evaluated in this study, it generally exhibited longer execution times, higher latest-value query latencies and greater storage requirements than the other database systems. Although TimescaleDB remained operationally stable throughout the experiments, its performance profile reflected different trade-offs compared with MongoDB and InfluxDB. The results indicate that extending a relational database architecture with time-series capabilities may involve additional overheads that affect ingestion performance, query responsiveness and storage utilization for the workload considered in this research.


The study also demonstrated that performance evaluation of time-series databases should not rely on a single metric. Databases that perform well in one area may exhibit different behavior when evaluated against other workload requirements. The results showed that ingestion performance, query latency and storage efficiency each provide a different perspective on database suitability. Consequently, database selection decisions should be guided by the specific operational requirements of the target application rather than by a single measure of performance.


Overall, the findings indicate that all three database systems were capable of supporting Smart Ocean environmental observation data, but they did so with different performance characteristics and trade-offs. Within the scope of the benchmark implementation and the performance metrics evaluated, MongoDB emerged as the strongest overall performer, InfluxDB demonstrated balanced and competitive behavior, and TimescaleDB provided functional support for the workload while exhibiting higher storage and latency costs. The study therefore concludes that database architecture significantly influences the efficiency with which large-scale time-series workloads are managed and that careful evaluation of workload requirements is essential when selecting a database platform for environmental monitoring systems such as Smart Ocean.


\section{Contributions of the Study}
This study contributes to the growing body of knowledge on the management of large-scale time-series data by providing both a practical benchmarking framework and empirical evidence regarding the performance characteristics of selected database technologies under environmental monitoring workloads. The contributions of the research span methodological, technical and practical dimensions



The first contribution is the design and implementation of a reproducible benchmarking framework for evaluating time-series database performance. The framework provided a standardized environment for executing ingestion benchmarks, collecting performance measurements and storing benchmark results. Through the use of a common execution workflow, database abstraction mechanisms and containerized deployment, the framework enabled MongoDB, InfluxDB and TimescaleDB to be evaluated under equivalent conditions. This implementation provides a foundation that can be extended for future benchmarking studies involving additional databases, larger datasets or alternative workload patterns.



The second contribution is the development of a representative benchmark workload based on Smart Ocean environmental observation data. Synthetic datasets were generated using a common observation model that captured the temporal, spatial and sensor characteristics of marine monitoring systems. By creating datasets of progressively increasing size and transforming them into formats appropriate for each database system, the study established a controlled mechanism for assessing how database performance changes with increasing data volume. The resulting workload provides a practical approximation of the challenges associated with storing and managing environmental sensor observations.



The third contribution is the generation of empirical performance data comparing MongoDB, InfluxDB and TimescaleDB under identical experimental conditions. The study produced quantitative evidence relating to ingestion performance, execution time, observation throughput, data throughput, storage efficiency and latest-value query latency. By evaluating all three database systems using the same datasets, deployment environment and benchmark methodology, the research provides a consistent basis for understanding the performance trade-offs associated with different database architectures.



The fourth contribution is the provision of evidence that can support database selection for large-scale time-series applications. The benchmark results demonstrate how document-oriented, relational time-series and purpose-built time-series database architectures behave when managing environmental observation data. The findings offer practical insights for researchers, engineers and system designers who must select database technologies for applications involving continuous sensor data generation, environmental monitoring and similar time-series workloads.



Finally, the study contributes to the Smart Ocean research initiative by providing an initial performance evaluation of database technologies capable of supporting underwater sensor network data management. The findings establish a baseline for future investigations into scalable storage solutions for environmental monitoring systems and provide a foundation for further research into the management of large-scale time-series data within Smart Ocean and related domains.


\section{Limitations of the Study}
While the study achieved its objectives and generated useful insights into the performance characteristics of MongoDB, InfluxDB and TimescaleDB, several limitations should be acknowledged when interpreting the results.


First, the evaluation was limited to three database technologies. Although MongoDB, InfluxDB and TimescaleDB represent widely used approaches to time-series data management, there are other database systems that were not included in the study. Consequently, the findings should be interpreted as a comparison among the selected technologies rather than as a comprehensive evaluation of all available time-series database solutions.


Second, the benchmark focused on a limited set of workload types. The evaluation concentrated primarily on data ingestion performance, storage efficiency and latest-value query latency. Other workload characteristics that may be important in operational environments, such as complex analytical queries, time-range aggregations, concurrent access patterns, data compression, retention management and long-term archival performance, were not included within the scope of the benchmark. As a result, the findings reflect performance under the specific workloads evaluated rather than all possible usage scenarios.


Third, the experiments were conducted within a single deployment environment using a fixed hardware and software configuration. Although containerization helped ensure consistency and reproducibility across benchmark executions, performance characteristics may vary when databases are deployed on different hardware platforms, operating systems or distributed infrastructures. The results therefore represent the behavior of the evaluated systems under the experimental conditions used in this study.

Fourth, the benchmark datasets were generated from synthetic environmental observations designed to represent Smart Ocean workloads. While this approach ensured consistency, repeatability and controlled dataset sizes, synthetic data may not fully capture all characteristics of real-world sensor streams, including irregular sampling intervals, missing values, data anomalies and varying observation patterns.



Finally, resource utilization metrics such as CPU usage, memory consumption, disk I/O activity and network utilization were not included in the final benchmark scope. The study therefore evaluated database performance primarily through execution time, throughput, query latency and storage efficiency. Additional insight into the operational costs of each database system could be obtained by incorporating resource utilization measurements into future evaluations.


Despite these limitations, the study provides a consistent and reproducible comparison of the selected database systems and offers valuable evidence regarding their performance characteristics under representative environmental monitoring workloads.



\section{Recommendations for Future Work}
The findings of this study provide a foundation for further research into the evaluation and selection of database technologies for large-scale time-series data management. While the benchmarking framework successfully compared MongoDB, InfluxDB and TimescaleDB using representative Smart Ocean workloads, several opportunities exist to extend and enhance the scope of future investigations.


One area for future work is the evaluation of additional database technologies. The current study focused on three widely used database systems representing document-oriented, purpose-built time-series and relational time-series architectures. Future studies could expand the comparison to include other platforms such as QuestDB, ClickHouse, Apache Druid, VictoriaMetrics or other emerging time-series data management solutions. Including a broader range of technologies would provide a more comprehensive understanding of the database landscape and offer additional guidance for practitioners selecting storage solutions for environmental monitoring systems.


Future research should also incorporate a wider variety of workload types. The benchmark developed in this study focused primarily on ingestion performance, storage efficiency and latest-value query latency. While these are important aspects of Smart Ocean data management, operational systems frequently perform more complex analytical tasks. Additional benchmark workloads could include time-range queries, aggregation operations, moving averages, statistical summaries, trend analysis and multi-parameter retrieval scenarios. Evaluating these workloads would provide a more complete picture of database performance across the full lifecycle of time-series data processing.


Another important direction is the investigation of distributed and clustered deployments. The benchmark experiments were conducted using single-node database instances to ensure consistency and simplify performance comparisons. However, large-scale environmental monitoring systems may eventually require distributed architectures to support higher ingestion rates, larger storage capacities and increased availability. Future studies could evaluate how database performance changes when deployed across multiple nodes and assess scalability, fault tolerance and load distribution characteristics under distributed operating conditions.

Resource utilization should also be incorporated into future benchmarking efforts. The current study evaluated performance using execution time, throughput, query latency and storage efficiency. While these metrics provide valuable insight into database behavior, they do not capture the computational resources required to achieve that performance. Future work should include measurements of CPU utilization, memory consumption, disk I/O activity and network usage. Such measurements would enable a more complete assessment of operational efficiency and help identify trade-offs between performance and infrastructure requirements.


Future studies could further investigate database behavior under concurrent workloads. Real-world environmental monitoring systems often support multiple users, applications and services accessing data simultaneously while new observations continue to be ingested. Benchmarking under concurrent read and write workloads would provide insight into how database performance changes under realistic operational conditions and would help evaluate system responsiveness in multi-user environments.



The benchmark datasets used in this study were generated synthetically to ensure repeatability and control over workload size. Future research could complement these experiments with real-world Smart Ocean datasets collected from operational sensor deployments. Evaluating database performance using real environmental observations would help determine whether the performance characteristics observed in this study remain consistent under more complex and potentially irregular data generation patterns.


Finally, the benchmarking framework developed in this research could be extended into a more comprehensive evaluation platform. Additional benchmark modules could be added to assess data compression effectiveness, retention policy management, backup and recovery performance, long-term archival storage behavior and advanced analytical capabilities. Such extensions would broaden the applicability of the framework and support more comprehensive evaluations of database technologies for environmental monitoring and other large-scale time-series applications.


By pursuing these research directions, future studies can build upon the foundation established in this thesis and develop a deeper understanding of how database architectures perform under the diverse requirements of modern time-series data management systems. Collectively, these efforts would contribute to more informed database selection decisions and support the continued development of scalable data management solutions for Smart Ocean and similar environmental monitoring initiatives.

% \section{Main Contributions}

% \section{Conclusions on Research Questions}

% \section{Threats to Validity}

% \section{Future Work}